%%%%%%%%%%%%%%%%%%%%%%%%%%%%%%%%%%%%%%%%%%%%%%%%%%%%%%%%%%%%%
%                                                           %
% JUAN LUIS VARONA - MARIO PEREZ RIERA                      %
%                                                           %
% BEGINNING DATE: 5 - MAY - 1998                            %
% LAST MODIFICATION DATE: 8 - MARZO - 1999                  %
%                                                           %
% Written with AMS-LaTeX 1.2                                %
%                                                           %
%%%%%%%%%%%%%%%%%%%%%%%%%%%%%%%%%%%%%%%%%%%%%%%%%%%%%%%%%%%%%


%%%%%%%%%%%%%%%%%%%%%%%%%%%%%%%%%%%%%%%%%%%%%%%%%%%%%%%%%%%%%
% Beginning of the paper
%%%%%%%%%%%%%%%%%%%%%%%%%%%%%%%%%%%%%%%%%%%%%%%%%%%%%%%%%%%%%
%
% AMS-LaTeX 1.2 file for journals:
%
\documentclass[reqno]{amsart}
% QUITAR LA LINEA ANTERIOR SI SE USA EL FORMATO LaTeX2e/amsart

%%%% \renewcommand{\baselinestretch}{2} % DOBLE ESPACIO
%%%% \setcounter{page}{3}

%\date{November 12, 1998.}

\date{PUBLISHED IN \textit{J.\ Math.\ Anal.\ Appl.}
\textbf{236} (1999), no.~1, 125--147.}

%%%%%%%%%%%%%%%%%%%%%%%%%%%%%%%%%%%%%%%%%%%%%%%%%%%%%%%%%%%%%%%%%%%%%%%

\theoremstyle{plain}
\newtheorem{Thm}{Theorem}
\newtheorem{Cor}[Thm]{Corollary} % Same numering sequence than Thm
\newtheorem{Lem}{Lemma}
\theoremstyle{definition}
\newtheorem{Defn}{Definition}
\newtheorem*{Ackn}{Acknowledgement}
\theoremstyle{remark}
\newtheorem*{Rem}{Remark}

\newcommand{\supp}{\operatorname{supp}} % support
\newcommand{\lpnorm}[3]{\|{#1}\|_{L^{#2}({#3})}} % L^p-norm of a function
\newcommand{\Id}{\operatorname{Id}} % identity map
\newcommand{\spa}{\operatorname{span}} % span 
\newcommand{\htr}{\mathcal{H}_\alpha} % Hankel transform of order \alpha
\newcommand{\hf}{{\,{}_2F_1}} % hypergeometric function

%%%%%%%%%%%%%%%%%%%%%%%%%%%%%%%%%%%%%%%%%%%%%%%%%%%%%%%%%%%%%%%%%%%%%%%%
% This is a fix suggested by the AMS TeX office to prevent overfulls in
% running heads containing accented characters
%
\makeatletter
\renewcommand{\leftmark}{%
  \vbox to\headheight{\vss
    \hbox{\expandafter \@firstoftwo \topmark {}{}}%
  }%
}
\makeatother
%%%%%%%%%%%%%%%%%%%%%%%%%%%%%%%%%%% End of the fix %%%%%%%%%%%%%%%%%%%%%

\begin{document}

\title[Fourier-Neumann series]{Mean and almost everywhere
convergence of Fourier-Neumann series}

\author[\'O. Ciaurri]{\'Oscar Ciaurri}
\address{Departamento de Matem\'aticas y Computaci\'on,
Universidad de La Rioja, Edificio J.~L.~Vives,
Calle Luis de Ulloa s/n, 26004 Logro\~no, Spain}
\email{osciaurr@dmc.unirioja.es}
\thanks{Research of the first author supported by grant ATUR97/013 of the UR}

\author[J. J. Guadalupe]{Jos\'e J. Guadalupe}
\address{Departamento de Matem\'aticas y Computaci\'on,
Universidad de La Rioja, Edificio J.~L.~Vives,
Calle Luis de Ulloa s/n, 26004 Logro\~no, Spain}
\email{joguadal@dmc.unirioja.es}

\author[M. P\'erez]{Mario P\'erez}
\address{Departamento de Matem\'aticas,
Universidad de Zaragoza, Edificio de Matem\'aticas,
Ciudad Universitaria s/n, 50009 Zaragoza, Spain}
\email{mperez@posta.unizar.es}

\author[J. L. Varona]{Juan L. Varona}
\address{Departamento de Matem\'aticas y Computaci\'on,
Universidad de La Rioja, Edificio J.~L.~Vives,
Calle Luis de Ulloa s/n, 26004 Logro\~no, Spain}
\email{jvarona@dmc.unirioja.es}
\urladdr{http://www.unirioja.es/dptos/dmc/jvarona/welcome.html}
\thanks{Research of the second, third and fourth authors supported by grants
PB96-0120-C03-02 of the DGES and API-98/B12 of the UR}

\subjclass{Primary 42C10; Secondary 44A05}

\keywords{Bessel functions, Fourier series, Neumann series, mean convergence,
almost everywhere convergence, Hankel transform.}

\begin{abstract}
Let $J_\mu$ denote the Bessel function of order $\mu$. The functions
$x^{-\alpha/2-\beta/2-1/2}J_{\alpha+\beta+2n+1}(x^{1/2})$, $n = 0, 1, 2, \dots$,
form an orthogonal system in $L^2((0,\infty),x^{\alpha+\beta}dx)$
when $\alpha+\beta>-1$.
In this paper we analyze the range of $p$, $\alpha$ and $\beta$ for which the
Fourier series with respect to this system converges in the
$L^p((0,\infty),x^\alpha dx)$-norm. Also, we describe the space in which the
$\spa$ of the system is dense and we show some of its properties.
Finally, we study the almost everywhere convergence of the Fourier series
for functions in such spaces.
\end{abstract}

%%%%%%%%%%%%%%%%%%%%%%%%%%%%%%%%%%%%%%%%%%%%%%%%%%%%%%%%%%%%%%%%%%%%%%%
\maketitle
%%%%%%%%%%%%%%%%%%%%%%%%%%%%%%%%%%%%%%%%%%%%%%%%%%%%%%%%%%%%%%%%%%%%%%%

%%%%%%%%%%%%%%%%%%%%%%%%%%%%%%%%%%%%%%%%%%%%%%%%%%%%%%%%%%%%%%%%%%%%%%%
\section{Introduction and notation}
%%%%%%%%%%%%%%%%%%%%%%%%%%%%%%%%%%%%%%%%%%%%%%%%%%%%%%%%%%%%%%%%%%%%%%%

Let $J_\mu$ be the Bessel function of order $\mu$.
For $\alpha>-1$, the formula
\begin{equation*}
   \int_0^\infty J_{\alpha+2n+1}(x) J_{\alpha+2m+1}(x) \, \frac{dx}{x}
   = \frac{\delta_{nm}}{2(2n+\alpha+1)},
   \quad n, m = 0,1,2, \dots
\end{equation*}
(see~\cite[Ch.~XIII, 13.41 (7), p.~404]{Wat} and~\cite[Ch.~XIII, 13.42 (1),
p.~405]{Wat}), provides an orthonormal system $\{j_n^\alpha\}_{n=0}^\infty$ in
$L^2((0,\infty), x^\alpha \, dx)$ [$L^2(x^\alpha)$, from now on], given by
\begin{equation*}
   j_n^\alpha(x)
   = \sqrt{\alpha+2n+1} J_{\alpha+2n+1}(\sqrt x) \, x^{-\alpha/2-1/2},
   \quad n = 0,1,2, \dots.
\end{equation*}
For each suitable function $f$, let $S_nf$ be the
$n$-th partial sum of its Fourier series with respect to the system
$\{j_n^\alpha\}_{n=0}^\infty$. Series of this kind are a particular case of
series $\sum_{n\ge0} a_n J_{\alpha+n}$, which are usually called Neumann
series, so that we will refer to $S_nf$ as a Fourier-Neumann series.
In~\cite{Var}, one of the authors studied the mean convergence in
$L^p(x^\alpha)$ of these Fourier series. In this context, some operators and
spaces were introduced. In this paper we extend these results and also
study the almost everywhere convergence.

For $\alpha > -1$, let us define the integral operator $\htr$ by
\begin{equation*}
    \htr(f,x) = \frac{x^{-\alpha/2}}{2}
     \int_0^\infty {f(t)J_\alpha (\sqrt {xt})t^{\alpha /2}\,dt},
     \quad x > 0,
\end{equation*}
for suitable functions $f$. This is a modified Hankel transform: the
(non-modified) Hankel transform is the integral operator with kernel
$J_\alpha(xt) (xt)^{1/2}$ and unweighted Lebesgue measure.
See~\cite{Dur},~\cite{St-Tr},~\cite{Me-So} for some modified and non-modified
Hankel transforms. In the case $\alpha \ge -1/2$, the Hankel transform
satisfies
\begin{equation*}
    \lpnorm{\htr f}\infty{x^\alpha}
      \le C \lpnorm{f}1{x^\alpha},
      \quad f \in L^1(x^\alpha)
\end{equation*}
with some constant $C$ independent of $f$.
Moreover, $\htr$ can be defined in $L^2(x^\alpha)$ satisfying
$\int_0^\infty (\htr f)g x^\alpha\,dx = \int_0^\infty (\htr g)f x^\alpha\,dx$,
$\htr^2 = \Id$, and $\lpnorm{\htr f}2{x^\alpha} =
\lpnorm{f}2{x^\alpha}$. From these results and interpolation we obtain
\begin{equation*}
    \lpnorm{\htr f}q{x^\alpha}
      \le C \lpnorm{f}p{x^\alpha},
      \quad f \in L^p(x^\alpha)
\end{equation*}
for $1\le p\le2$, where $q$ denotes, here and in the rest of the paper, the
conjugate of $p$, that is, $1/p+1/q=1$.

The Hankel transform of the function $j_n^\alpha$ is
\begin{equation*}
   \htr(j_n^\alpha,x) = \sqrt{\alpha+2n+1} P_n^{(\alpha,0)}(1-2x)
   \chi_{[0,1]}(x),
\end{equation*}
where $P_n^{(\alpha,\beta)}(x)$ is the $n$-th Jacobi polynomial of order
$(\alpha,\beta)$; see, for instance,~\cite[Ch.~8.11, (5), p.~47]{E-TIT} (a
thorough description of Jacobi polynomials can be found in~\cite[Ch.~X]{E-HTF}
or~\cite{Sze}). 

\begin{Rem}
There is a delicacy with this formula. Actually, $\htr$ was defined, as a
first step, as a Lebesgue integral for suitably integrable functions. Then,
$\htr$ is extended to $L^p$ spaces where the integral representation is no
longer valid for some functions. Now, the integrals from~\cite[Ch.~8.11, (5),
p.~47]{E-TIT} are improper Riemann integrals. Hence, the proper understanding
of those integrals should be
\begin{equation*}
   \lim_{N\to\infty}
   \frac{x^{-\alpha/2}}{2}
   \int_0^N {j_n^\alpha(t) J_\alpha (\sqrt {xt})t^{\alpha /2}\,dt} =
   \sqrt{\alpha+2n+1} P_n^{(\alpha,0)}(1-2x) \chi_{[0,1]}(x).
\end{equation*}
Since $j_n^\alpha \chi_{[0,N]}$ is an integrable function, the integral form
of $\htr$ is valid here and we can conclude that
\begin{equation*}
   \lim_{N\to\infty}
   \htr(j_n^\alpha \chi_{[0,N]},x)
   = \sqrt{\alpha+2n+1} P_n^{(\alpha,0)}(1-2x) \chi_{[0,1]}(x),
\end{equation*}
where the limit holds in the almost everywhere sense. Finally, the $L^p$
boundedness of the operator $\htr$ for $\frac{4(\alpha+1)}{2\alpha+3} < p < 2$
and the fact that $\lim_{N\to\infty} j_n^\alpha \chi_{[0,N]} = j_n^\alpha$ in
$L^p$ yields $\htr(j_n^\alpha,x) = \sqrt{\alpha+2n+1} P_n^{(\alpha,0)}(1-2x)
\chi_{[0,1]}(x)$ in $L^p$. Similar comments apply to Lemma~\ref{lem:HTjnab}
below.
\end{Rem}

Since the Hankel transform of $j_n^\alpha$ is supported on
$[0,1]$, not every function $f \in L^p(x^\alpha)$, $1 < p \leq 2$,
can be approximated in norm by its Fourier series $S_nf$. As a first approach,
any such function should, at least, have its Hankel transform supported
on $[0,1]$. But we also deal with spaces $L^p(x^\alpha)$, $p>2$,
where $\htr$ is not defined and so, we need to describe the functions that
we want to approximate in a different, but, in some sense, similar way.

The main tool here is $M_\alpha$, the multiplier for the Hankel transform. For
$\alpha \geq -1/2$ and $-\frac14 < (\alpha+1)(\frac12-\frac1p) < \frac14$,
$M_\alpha$ is a bounded operator from $L^p(x^\alpha)$ into itself (this is
known as Herz's theorem, see~\cite{Herz}). Also,
\begin{equation*}
    \htr(M_\alpha f) = \htr(f) \chi_{[0,1]}
\end{equation*}
for $f \in L^p(x^\alpha) \cap L^2(x^\alpha)$, $M_\alpha^2 f = M_\alpha f$ for
$f \in L^p(x^\alpha)$ and
\begin{equation}
\label{ec:Mafg}
    \int_0^\infty f(x) M_\alpha(g,x) x^\alpha\,dx
    = \int_0^\infty g(x) M_\alpha(f,x) x^\alpha\,dx
\end{equation}
for $f \in L^p(x^\alpha)$ and $g \in L^q(x^\alpha)$.

\begin{Defn}
For each $\alpha$ and $p$ with $\alpha \geq -1/2$ and
$-\frac14 < (\alpha+1)(\frac12-\frac1p) < \frac14$, let us define the
$L^p(x^\alpha)$ subspace
\begin{equation*}
    E_{p,\alpha} = \{ f \in L^p(x^\alpha): M_\alpha f = f\}
    = M_\alpha(L^p(x^\alpha)).
\end{equation*}
\end{Defn}

It is clear that, for $f \in E_{p,\alpha} \cap L^2(x^\alpha)$, the Hankel
transform of $f$ is supported on $[0,1]$ and so these spaces are suitable
for our purposes.
The spaces $E_{p,\alpha}$ have some interesting properties: For $s<r$,
$E_{s,\alpha} \subset E_{r,\alpha}$ and the inclusion is continuous and dense.
Besides, the dual space is $(E_{p,\alpha})' = E_{q,\alpha}$.

Let us also consider, for each $\alpha \geq -1$ and each suitable $p$ (we will
go into the details later), the $L^p(x^\alpha)$ subspace
\begin{equation*}
	B_{p,\alpha} = \overline{\spa} \{j_n^\alpha(x)\}_{n=0}^\infty
   \quad (\text{closure in } L^p(x^\alpha)).
\end{equation*}
In~\cite{Var}, one of us showed that $S_nf \to f$ in the $L^p(x^\alpha)$-norm
for any $f \in B_{p,\alpha}$, if $\alpha \geq -1/2$ and
\begin{equation*}
    \max\left\{ \frac43, \frac{4(\alpha+1)}{2\alpha+3} \right\}
    < p <
    \min\left\{ 4, \frac{4(\alpha+1)}{2\alpha+1}\right\};
\end{equation*}
moreover, for this range of $p$, we showed that
$B_{p,\alpha} = E_{p,\alpha}$. Therefore,
$\{j_n^\alpha\}_{n=0}^\infty$ is a basis for the space $E_{p,\alpha}$. By the
way, notice that for $\alpha \geq -1/2$,
\begin{gather*}
   \frac{4(\alpha+1)}{2\alpha+3} < p
   \iff -\frac14 < (\alpha+1) \left( \frac12 - \frac1p \right), \\
   p < \frac{4(\alpha+1)}{2\alpha+1}
   \iff (\alpha+1) \left( \frac12 - \frac1p \right) < \frac 14.
\end{gather*}

Our purpose in this paper is to improve and extend these convergence
results, and show additional properties of the $E_{p,\alpha}$ spaces.
In particular, we find some conditions on $\alpha$, $\beta$, and $p$ under
which the functions $j_n^{\alpha+\beta}$ ($n=0,1,2,\dots$) are a basis for
$E_{p,\alpha}$. For instance, $\beta$ can be taken so that $\alpha + \beta$ is
half an integer, which makes the functions $j_n^{\alpha+\beta}$ better known.
The almost everywhere convergence of $S_nf$ is studied, as well.

Also, we can interpret the convergence in the following way:
changing the parameters, we take
$\{j_n^{\lambda}\}_{n=0}^\infty$, which is orthogonal in $L^2(x^\lambda)$,
and we study the convergence in $L^p(x^\mu)$. This is a typical situation
in the study of mean convergence of Fourier series. For instance, in the case
of Jacobi polynomials $\{P_n^{(\alpha,\beta)}(x)\}_{n=0}^\infty$,
Pollard~\cite{Pol-III} studied the convergence in the natural space
$L^p((-1,1),(1-x)^\alpha(1+x)^\beta)$ and, later, Muckenhoupt~\cite{Muck}
described the behaviour in $L^p((-1,1),(1-x)^a(1+x)^b)$. Similar situations
occur with other orthogonal systems (Laguerre, Hermite, Freud weights,
Bessel and Dini).

We are interested in the approximation of functions in $L^p(x^\alpha)$
by Fourier series in the system $\{j_n^{\alpha+\beta}\}_{n=0}^\infty$.
So, our first target is to determine the range of $p$, $\alpha$, and $\beta$
for which $j_n^{\alpha+\beta} \in L^p(x^\alpha)$ for all $n \in \mathbb{N}$.
We do this in section~\ref{sec:Bpab}.

In section~\ref{sec:convergence} we state some of the main results of this
paper: the uniform boundedness and convergence of the partial sum operator of
Fourier-Neumann series. The proofs are given in sections~\ref{sec:proofboundSn}
and~\ref{sec:proofconvae}. The mean convergence can only hold for functions in
the closure of the linear combinations of the functions $j_n^{\alpha+\beta}$.
In section~\ref{sec:Bpab=Epa} this space is shown to coincide with
$E_{p,\alpha}$ under some conditions on $p$, $\alpha$, and $\beta$. Some
applications are given in section~\ref{sec:applications}.

Throughout this paper, unless otherwise stated, we will use $C$ (or $C_1$)
to denote a positive constant independent of $n$ (and all other variables),
which can assume different values in different occurrences. Also, in what
follows, $a_n \sim b_n$, for $a_n, b_n>0$, means $C \le a_n/b_n \le C_1$.

%%%%%%%%%%%%%%%%%%%%%%%%%%%%%%%%%%%%%%%%%%%%%%%%%%%%%%%%%%%%%%%%%%%%%%%
\section{The spaces $B_{p,\alpha,\beta}$}
\label{sec:Bpab}
%%%%%%%%%%%%%%%%%%%%%%%%%%%%%%%%%%%%%%%%%%%%%%%%%%%%%%%%%%%%%%%%%%%%%%%

We will use here the well known estimates (see~\cite{E-HTF}
or~\cite[Ch.~III, 3.1 (8), p.~40]{Wat} and~\cite[Ch.~VII, 7.21 (1),
p.~199]{Wat})
\begin{equation}
\label{ec:bessel0}
   J_\mu(x) = \frac{x^\mu}{2^\mu \Gamma(\mu+1)} + O(x^{\mu+2}),
   \quad x\to 0+,
\end{equation}
and
\begin{equation}
\label{ec:besselinfty}
   J_\mu(x) = \sqrt{\frac2{\pi x}}
   \left[\cos\left(x-\frac{\mu\pi}2-\frac\pi4\right) + O(x^{-1})\right],
   \quad x\to\infty.
\end{equation}

\begin{Lem}
\label{lem:normajn}
Let $a>-1$, $1 < p < \infty$. Then, $j_n^a \in L^p(x^b)$ for all
$n=0,1,2,\dots$ if and only if $b > -1$ and $-\frac14 < (b+1)\left(\frac12 -
\frac1p\right) + \frac{a-b}2$. Furthermore, in this case,
\begin{equation*}
   \|j_n^a\|_{L^p(x^b)} \sim
      \begin{cases}
        n^{-(a+1) + 2(b+1)/p},
        &\text{ if } p < 4, \\[2pt]
        n^{-(2a-b+1)/2} (\log n)^{1/4},
        &\text{ if } p = 4, \\[2pt]
        n^{-(5/6+a) + (6b + 4)/(3p)},
        &\text{ if } p > 4.
      \end{cases}
\end{equation*}
\end{Lem}

\begin{proof}
Inequalities $b > -1$ and $-\frac14 < (b+1)\left(\frac12 - \frac1p\right) +
\frac{a-b}2$ follow from~(\ref{ec:bessel0}) and~(\ref{ec:besselinfty}). Then,
estimates such as~(\ref{ec:jnu}) below (see~\cite{Bar-Cor} and~\cite{Cord}) 
show that $\|j_n^a\|_{L^p(x^b)}$ is bounded above by a constant times the
right hand side. The lower bound follows from more precise estimates for the
Bessel functions, as shown in~\cite{Bar-Cor} and~\cite{Cord}. For a similar
expression, see~\cite{Stem}.
\end{proof}

As a consequence, the following definition makes sense.

\begin{Defn}
\label{defn:Bp}
For each $\alpha$, $\beta$, and $p$ with $\alpha > -1$, $\alpha + \beta > -1$,
$1< p < \infty$, and
\begin{equation*}
    -\frac14 < (\alpha+1)\left(\frac12-\frac1p\right) + \frac{\beta}2,
\end{equation*}
let us define
\begin{equation*}
	  B_{p,\alpha,\beta} = \overline{\spa} \{j_n^{\alpha+\beta}(x)\}_{n=0}^\infty
   \quad (\text{closure in } L^p(x^\alpha)).
\end{equation*}
\end{Defn}

Note that we assume $\alpha>-1$ in the definition of $B_{p,\alpha,\beta}$;
however, we require $\alpha\ge-1/2$ for $E_{p,\alpha}$. Actually, the
boundedness of $M_\alpha$ can be studied also for $\alpha>-1$, so that the
definition of $E_{p,\alpha}$ can be extended to the whole range $\alpha >
-1$. But in the case $\alpha < -1/2$, the $\htr$ transform does not have as
good properties as in the case $\alpha\ge-1/2$. As a consequence, the spaces
$E_{p,\alpha}$ do not behave for $\alpha<-1/2$ like for $\alpha\ge-1/2$. Thus,
some of the results in this paper will be established for $\alpha>-1$, but we
will require $\alpha\ge-1/2$ when $E_{p,\alpha}$ appears.

The following Lemma proves that $B_{p,\alpha} \subset B_{p,\alpha,\beta}$,
under some conditions on $\alpha$, $\beta$, and $p$.

\begin{Lem}
\label{lem:jnajna+b}
Let $\alpha>-1$, $\alpha + \beta > -1$, and $1<p<4$ such that
\begin{equation*}
    \frac{-1}4 < (\alpha+1)\left(\frac12-\frac1p\right) - \frac{|\beta|}2.
\end{equation*}
Then,
\begin{equation}
\label{ec:jnajna+b}
    j_n^\alpha = \sum_{k=n}^\infty a_{n,k} j_k^{\alpha+\beta}
\end{equation}
pointwise and in $L^p(x^\alpha)$, where
\begin{equation}
\label{ec:ank}
    a_{n,k} = \frac{2^\beta \, \sqrt{\alpha + 2n + 1} \,
        \sqrt{\alpha+\beta+2k+1} \,
        \Gamma(1-\beta) \, \Gamma(\alpha+\beta+k+n+1)}
       {\Gamma(1+k-n) \, \Gamma(1-\beta-k+n)\,
        \Gamma(\alpha+k+n+2)}.
\end{equation}
\end{Lem}

\begin{Rem}
If $\beta \in \mathbb{N}$, then $\Gamma(1-\beta) / \Gamma(1-\beta-k+n)$ should
be replaced by $-\beta (-\beta -1) (-\beta -2) \dots (1-\beta-k+n)$ in
formula~(\ref{ec:ank}).
\end{Rem}

\begin{proof}
The pointwise convergence and (\ref{ec:ank}) follow from~\cite[Ch.~V,
5.21 (1), p.~139]{Wat} (conditions $\alpha > -1$ and $\alpha + \beta > -1$ are
required). Strictly speaking, condition $\beta \notin \mathbb{N}$ should also
be assumed, following~\cite{Wat}. But this is only a formal requirement to get
$a_{n,k}$ in the form of~(\ref{ec:ank}).

For the $L^p$ convergence, we need only prove that the series converges:
that the sum is precisely $j_n^\alpha$ will then follow from the fact that
this holds in the almost everywhere sense.

If $\beta$ is an integer, then there are only finitely many $a_{n,k} \not= 0$
and the series in~(\ref{ec:jnajna+b}) is a finite sum. If $\beta$ is not an
integer, Stirling's formula for the Gamma function gives, for each fixed
$n$,
\begin{equation*}
    |a_{n,k}| \sim k^{2\beta-3/2}, \quad k \to \infty.
\end{equation*}
Also, from Lemma~\ref{lem:normajn}, $p < 4$ and $\frac{-1}4 <
(\alpha+1)\left(\frac12-\frac1p\right) + \frac{\beta}2$, we have
\begin{equation*}
    \lpnorm{j_k^{\alpha+\beta}}{p}{x^\alpha}
      \sim k^{-(\alpha+\beta+1)+2(\alpha+1)/p}.
\end{equation*}
These estimates and $\frac{-1}4 < (\alpha+1)\left(\frac12-\frac1p\right) -
\frac{\beta}2$ prove that
\begin{equation*}
    \sum_{k=n}^\infty |a_{n,k}| \, \lpnorm{j_k^{\alpha+\beta}}{p}{x^\alpha}
   < \infty.
\end{equation*}
\end{proof}

%%%%%%%%%%%%%%%%%%%%%%%%%%%%%%%%%%%%%%%%%%%%%%%%%%%%%%%%%%%%%%%%%%%%%%%
\section{Uniform boundedness and convergence of Fourier-Neumann series}
\label{sec:convergence}
%%%%%%%%%%%%%%%%%%%%%%%%%%%%%%%%%%%%%%%%%%%%%%%%%%%%%%%%%%%%%%%%%%%%%%%

Let us consider the partial sums of the Fourier series with respect to the
system
$\{j_n^{\alpha+\beta}\}_{n=0}^\infty$:
\begin{equation*}
   S_n(f,x) = \sum_{k=0}^n {c_k(f) j_k^{\alpha+\beta} (x)}, \qquad
   c_k(f) = \int_0^\infty {f(t) j_k^{\alpha+\beta}(t) t^{\alpha+\beta}\,dt}.
\end{equation*}
We are interested in the study of the uniform boundedness of
the partial sum operators
\begin{equation*}
    S_n \colon L^p(x^\alpha) \longrightarrow L^p(x^\alpha).
\end{equation*}
Our result is the following:

\begin{Thm}
\label{th:boundSn}
Let $\alpha>-1$, $\alpha+\beta>-1$, and $1 < p < \infty$. There exists a
constant $C>0$ such that
\begin{equation*}
    \lpnorm{S_nf}p{x^\alpha} \le C\lpnorm{f}p{x^\alpha},
    \quad f \in L^p(x^\alpha), \quad n \in \mathbb{N},
\end{equation*}
if and only if $\frac43<p<4$ and
\begin{equation}
\label{ec:range}
\begin{split}
   - \frac{\alpha+\beta+1}2
   &< (\alpha+1) \left( \frac12 - \frac1p \right) + \frac\beta2,
   \\
   - \frac14
   &< (\alpha+1) \left( \frac12 - \frac1p \right) + \frac\beta2 < \frac14.
\end{split}
\end{equation}
\end{Thm}

\begin{proof}
See section~\ref{sec:proofboundSn}.
\end{proof}

\begin{Cor}
\label{cor:conv}
Let $\alpha > -1$, $\alpha+\beta> -1$, $\frac43 < p < 4$, and
\begin{equation*}
\begin{split}
   - \frac{\alpha+\beta+1}2
   &< (\alpha+1) \left( \frac12 - \frac1p \right) + \frac\beta2,
   \\
   - \frac14
   &< (\alpha+1) \left( \frac12 - \frac1p \right) + \frac\beta2 < \frac14.
\end{split}
\end{equation*}
Then, $S_nf \to f$ in $L^p(x^\alpha)$ for all $f \in B_{p,\alpha,\beta}$.
\end{Cor}

\begin{proof}
$B_{p,\alpha,\beta}$ is the closure in $L^p(x^\alpha)$ of the orthogonal
system, so this is just a standard consequence of Theorem~\ref{th:boundSn}.
\end{proof}

Regarding the almost everywhere convergence of Fourier-Neumann series, we have
the following.

\begin{Thm}
\label{th:convae}
Let $\alpha > -1$, $\alpha+\beta> -1$, $\frac43 < p < 4$, and
\begin{equation*}
\begin{split}
   - \frac{\alpha+\beta+1}2
   &< (\alpha+1) \left( \frac12 - \frac1p \right) + \frac\beta2,
   \\
   - \frac14
   &< (\alpha+1) \left( \frac12 - \frac1p \right) + \frac\beta2 < \frac14.
\end{split}
\end{equation*}
Then, $S_nf \to f$ almost everywhere for any $f \in B_{p,\alpha,\beta}$.
\end{Thm}

\begin{proof}
See section~\ref{sec:proofconvae}.
\end{proof}


%%%%%%%%%%%%%%%%%%%%%%%%%%%%%%%%%%%%%%%%%%%%%%%%%%%%%%%%%%%%%%%%%%%%%%%
\section{The Hankel transform of order $\alpha$ for $j_n^{\alpha+\beta}$\\
and the spaces $B_{p,\alpha,\beta}$ and $E_{p,\alpha}$}
\label{sec:Bpab=Epa}
%%%%%%%%%%%%%%%%%%%%%%%%%%%%%%%%%%%%%%%%%%%%%%%%%%%%%%%%%%%%%%%%%%%%%%%

Theorem~\ref{th:boundSn} and Corollary~\ref{cor:conv} are more interesting if
we can describe the space $B_{p,\alpha,\beta}$. In this section, we will find
some conditions under which $B_{p,\alpha,\beta}$ and $E_{p,\alpha}$ coincide.

As we pointed out, the Hankel transform of order $\alpha$ for $j_n^\alpha$ can
be written in terms of the $n$-th Jacobi polynomial of order $(\alpha,0)$. It
is not difficult to obtain $\htr(j_n^{\alpha+\beta})$ from known results about
integrals of products of Bessel functions that can be expressed in terms of
hypergeometric $\hf$ functions. But the relation between
$\htr(j_n^{\alpha+\beta})$ and the Jacobi polynomials of order
$(\alpha,\beta)$ is not easily found in the literature. For instance, it does
not appear in the standard references~\cite{E-HTF,Wat, E-TIT}. For the sake of
completeness, in this section we obtain $\htr(j_n^{\alpha+\beta})$ explicitly
in terms of $P_n^{(\alpha,\beta)}$.

\begin{Lem}
\label{lem:HTjnab}
For $\alpha,\beta>-1$ with $\alpha + \beta > -1$,
\begin{equation*}
    \htr(j_n^{\alpha+\beta},x)
      = 2^{-\beta} \tfrac{\sqrt{\alpha+\beta+2n+1} \, \Gamma(n+1)}
        {\Gamma(\beta+n+1)}
        (1-x)^\beta P^{(\alpha,\beta)}_n(1-2x) \chi_{[0,1]}(x).
\end{equation*}
In particular,
$\supp (\htr(j^{\alpha+\beta}_n)) \subseteq [0,1]$.
\end{Lem}

\begin{proof}
We use the formula
\begin{multline}
\label{ec:hforig}
   \int_0^\infty t^{-\lambda} J_\mu(at) J_\nu(bt)\,dt
   \\
   = \frac{b^\nu a^{\lambda-\nu-1}
         \Gamma(\frac{\mu+\nu-\lambda+1}2)}
         {2^\lambda \Gamma(\nu+1) \Gamma(\frac{\lambda+\mu-\nu+1}2)}
       \hf\left( \frac{\mu+\nu-\lambda+1}2,
          \frac{\nu-\lambda-\mu+1}2; \nu+1;
          \frac{b^2}{a^2} \right),
\end{multline}
valid when $0<b<a$, $\mu+\nu-\lambda>-1$ and $\lambda>-1$; here, $\hf$ denotes
the hypergeometric function (see~\cite[Ch.~8.11, (9), p.~48]{E-TIT} or
\cite[Ch.~XIII, 13.4 (2), p.~401]{Wat}).

Taking $a=1$ and $x=b^2$ in~(\ref{ec:hforig}), and making the
corresponding changes of variable and parameters ($\nu=\alpha$,
$\mu=\alpha+\beta+2n+1$, $\lambda=\beta$) we get
\begin{equation*}
    \htr(j_n^{\alpha+\beta},x)
    = \tfrac{\sqrt{2n+\alpha+\beta+1}\,\Gamma(\alpha+n+1)}
      {2^\beta \Gamma(\beta+n+1) \Gamma(\alpha+1)}
      \hf(\alpha+n+1,-n-\beta;\alpha+1;x),
\end{equation*}
which is valid for $\alpha>-1$ and $\beta>-1$ in the interval $0<x<1$.
Now, we have
\begin{equation*}
    \hf(\alpha+n+1,-n-\beta;\alpha+1;x)
    = (1-x)^\beta \hf(-n,\alpha+\beta+n+1;\alpha+1;x),
\end{equation*}
where $\alpha,\beta>-1$, $n=0,1,2,\dots$, and
\begin{equation*}
    P_n^{(\alpha,\beta)}(x) =
      \tfrac{\Gamma(n+\alpha+1)}{\Gamma(\alpha+1)\Gamma(n+1)}
      \hf(-n, \alpha+\beta+n+1; \alpha+1; \tfrac{1-x}2),
      \quad \alpha, \beta > -1.
\end{equation*}
Therefore,
\begin{equation*}
    \htr(j_n^{\alpha+\beta},x)
      = \tfrac{\sqrt{\alpha+\beta+2n+1} \, \Gamma(n+1)}
        {2^\beta \Gamma(\beta+n+1)} (1-x)^\beta P^{(\alpha,\beta)}_n(1-2x),
        \quad x \in (0,1).
\end{equation*}
Now, let us calculate $\htr(j_n^{\alpha+\beta},x)$ for $x>1$. To do that, let
us take $x=a^2$, $b=1$, $\nu=\alpha+\beta+2n+1$, $\mu=\alpha$, and
$\lambda=\beta$ in~(\ref{ec:hforig}). In this way,
$\frac12(\lambda+\mu-\nu+1) = 0,-1, -2, \dots$, so the coefficient
$1/\Gamma(\frac12(\lambda+\mu-\nu+1))$ vanishes and we get
$\htr(j_n^{\alpha+\beta},x) = 0$.
\end{proof}

\begin{Thm}
\label{th:Bp=Ep}
Let $\alpha\ge-1/2$, $\beta > -1/2$, $4/3 < p$, with
\begin{equation*}
   -\frac14 < (\alpha+1) \left( \frac12 - \frac1p \right) < \frac14.
\end{equation*}
If $p < 2$, assume further
\begin{equation*}
   -\frac14 < (\alpha+1) \left(\frac12 - \frac1p \right) - \frac{|\beta|} 2. 
\end{equation*}
Then $B_{p,\alpha,\beta} = E_{p,\alpha}$.
\end{Thm}
   
\begin{proof}
{\it Case $p=2$.} The spaces $B_{2,\alpha,\beta}$ and $E_{2,\alpha}$
are well defined. Also $M_\alpha(j_n^{\alpha+\beta}) = j_n^{\alpha+\beta}$. In
other words, $B_{2,\alpha,\beta} \subset E_{2,\alpha}$. If they were not
equal, by the Hahn-Banach theorem there should exist some $T \in
(E_{2,\alpha})'$, $T \ne 0$, such that $T(j_n^{\alpha+\beta}) = 0$ $\forall
n$. But $(E_{2,\alpha})' = E_{2,\alpha}$, so there exists $\varphi \in
E_{2,\alpha}$, $\varphi\ne0$, such that $\int_0^\infty \varphi
j_n^{\alpha+\beta} x^\alpha \, dx = 0$ for every
$n$. Then
\begin{align*}
    0 &= \int_0^\infty \varphi j_n^{\alpha+\beta} x^\alpha \, dx
    = \int_0^\infty (\htr \varphi) (\htr j_n^{\alpha+\beta}) x^\alpha \, dx \\
    &= k_n \int_0^1 (\htr \varphi) P_n^{(\alpha,\beta)}(1-2x)
      (1-x)^\beta x^\alpha \, dx
\end{align*}
for every nonnegative integer $n$. Now, the Jacobi polynomials
$P_n^{(\alpha,\beta)}(x)$ are a complete orthogonal system with respect to the
measure $(1-x)^\alpha(1+x)^\beta \, dx$ on $(-1,1)$. A change of variable
proves that the polynomials $P_n^{(\alpha,\beta)}(1-2x)$ are a complete
orthogonal system with respect to the measure $(1-x)^\beta x^\alpha \, dx$ on
$(0,1)$. Thus, $\htr \varphi = 0$ on $(0,1)$. Since $\varphi \in
E_{2,\alpha}$, we also have $\htr \varphi = 0$ on $(1,\infty)$. Therefore,
$\htr \varphi = 0$ and we arrive at the contradiction $\varphi = 0$.

{\it Case $p>2$.} Note that $\alpha$, $\beta$, and $p$ meet the
requirements of Definition~\ref{defn:Bp}. Also, by the preceding case, we
have $j_n^{\alpha+\beta}\in E_{2,\alpha} \subset E_{p,\alpha}$. Thus,
$B_{p,\alpha,\beta}\subset E_{p,\alpha}$.

Now, let $f\in E_{p,\alpha}$. Given $\varepsilon>0$, there exists a function
$g\in L^2(x^\alpha)\cap L^p(x^\alpha)$ such that
$\lpnorm{f-g}{p}{x^\alpha}<\varepsilon$. Let $h=M_\alpha g$; then $h \in
L^2(x^\alpha)\cap L^p(x^\alpha)$ and $M_\alpha h = h$, so that $h\in
E_{2,\alpha}\cap E_{p,\alpha} = B_{2,\alpha,\beta}\cap E_{p,\alpha}$. Since
$M_\alpha$ is continuous, $\lpnorm{f-h}{p}{x^\alpha} = \lpnorm{M_\alpha
f-M_\alpha g}{p}{x^\alpha} < C\varepsilon$. As $h\in B_{2,\alpha,\beta}$,
there exists $h'\in \spa\{j_n^{\alpha+\beta}\}_{n=0}^\infty$ such that
$\lpnorm{h-h'}{2}{x^\alpha}<\varepsilon$. The inclusion $E_{2,\alpha} \subset
E_{p,\alpha}$ gives $\lpnorm{h-h'}{p}{x^\alpha}< C_1 \varepsilon$, so that, by
the triangle inequality, $\lpnorm{f-h'}{p}{x^\alpha}<C_2\varepsilon$.
This gives the inclusion $E_{p,\alpha}\subset B_{p,\alpha,\beta}$.

{\it Case $p<2$.} By Lemmas~\ref{lem:normajn} and~\ref{lem:HTjnab},
$j_n^{\alpha+\beta} \in L^2(x^\alpha)$ and $\htr(j_n^{\alpha+\beta})$ is
supported on $[0,1]$, so that $M_\alpha j_n^{\alpha+\beta} =
j_n^{\alpha+\beta}$. Since $j_n^{\alpha+\beta} \in L^p(x^\alpha)$ by
Lemma~\ref{lem:normajn}, it follows that $j_n^{\alpha+\beta} \in
E_{p,\alpha}$. Therefore, $B_{p,\alpha,\beta} \subset E_{p,\alpha}$.

The equality will follow if we prove that the only operator $T \in
(E_{p,\alpha})'$ such that $T(f) = 0$ for all $f \in B_{p,\alpha,\beta}$ is $T
= 0$. For such an operator, we have $T(j_n^\alpha) = 0$ for every $n \ge 0$,
since $j_n^\alpha \in B_{p,\alpha,\beta}$ by Lemma~\ref{lem:jnajna+b}. On the
other hand, by the duality $(E_{p,\alpha})'=E_{q,\alpha}$, where $1/p+1/q=1$,
there exists some $\varphi \in E_{q,\alpha}$ such that
\begin{equation*}
    T(f) = \int_0^\infty \varphi f x^\alpha\,dx, \quad f \in E_{p,\alpha}.
\end{equation*}
In particular,
\begin{equation}
\label{ec:cn=0}
   \int_0^\infty \varphi j_n^\alpha x^\alpha \, dx = 0, \quad n \ge 0.
\end{equation}
Under the present conditions on $p$ and $\alpha$, the preceding case
gives $B_{q,\alpha,0} = E_{q,\alpha}$, so that $\varphi \in B_{q,\alpha,0}$.
This, together with (\ref{ec:cn=0}) and Corollary~\ref{cor:conv}, gives
$\varphi = 0$.
\end{proof}

%%%%%%%%%%%%%%%%%%%%%%%%%%%%%%%%%%%%%%%%%%%%%%%%%%%%%%%%%%%%%%%%%%%%%%%
\section{Applications}
\label{sec:applications}
%%%%%%%%%%%%%%%%%%%%%%%%%%%%%%%%%%%%%%%%%%%%%%%%%%%%%%%%%%%%%%%%%%%%%%%

Some properties of the spaces $E_{p,\alpha}$ can be obtained from
Theorem~\ref{th:Bp=Ep}. Two examples will be given here, after this
preliminary result. 

\begin{Cor}
\label{cor:convMa}
Let $\alpha\ge-1/2$, $\beta=0$ and $4/3 < p < 4$ verifying
\begin{equation*}
   -\frac14 < (\alpha+1) \left( \frac12 - \frac1p \right) < \frac14.
\end{equation*}
Then, $S_nf \to M_\alpha f$ in $L^p(x^\alpha)$ and almost everywhere for all $f
\in L^p(x^\alpha)$.
\end{Cor}

\begin{proof}
Let $f \in L^p(x^\alpha)$, and so $M_\alpha f \in E_{p,\alpha}$.
Then, by Theorem~\ref{th:Bp=Ep} and Theorem~\ref{th:convae},
$S_n(M_\alpha f) \to M_\alpha f$ in $L^p(x^\alpha)$ and almost everywhere.
So, we only need to show that $S_n(M_\alpha f) = S_n(f)$, and
this is clear because, by~(\ref{ec:Mafg}),
\begin{equation*}
    \int_0^\infty (M_\alpha f) j_n^{\alpha+\beta} x^\alpha \,dx
    = \int_0^\infty f (M_\alpha j_n^{\alpha+\beta}) x^\alpha \,dx
    = \int_0^\infty f j_n^{\alpha+\beta} x^\alpha \,dx.
\end{equation*}
\end{proof}

\begin{Cor}
\label{cor:EpaEpa+b}
Let $\alpha\ge-1/2$, $-1/2 < \beta < 1$, $4/3 < p < 4$ with
$\alpha+\beta\ge-1/2$,
\begin{gather*}
    - \frac14 < (\alpha+\beta+1) \left(\frac12 - \frac1p \right) < \frac14,
   \\
   \max\left\{ -\frac14, -\frac14 - \frac\beta2, - \frac{\alpha+1}2 - \beta
       \right\}
   < (\alpha+1) \left(\frac12 - \frac1p \right) <
   \min\left\{ \frac14, \frac14 - \frac\beta2 \right\}.
\end{gather*}
If $p<2$, assume further $\beta<1/2$ and $-\frac14 + \frac\beta2 <
(\alpha+1) \left(\frac12 - \frac1p \right)$.

Then, $E_{p,\alpha} \cap L^p(x^{\alpha+\beta}) \subset E_{p,\alpha+\beta}$.
\end{Cor}

\begin{proof}
Let $f \in E_{p,\alpha} \cap L^p(x^{\alpha+\beta})$. By Theorem~\ref{th:Bp=Ep}
and Theorem~\ref{th:convae}, $S_n f \to f$ almost everywhere. Since $f \in
L^p(x^{\alpha+\beta})$, Corollary~\ref{cor:convMa} (with $\alpha +\beta$
instead of $\alpha$) gives $S_n f \to M_{\alpha+\beta}f$ almost everywhere.
Then, $f = M_{\alpha+\beta}f$ almost everywhere, that is,
$f \in E_{p,\alpha+\beta}$.
\end{proof}

\begin{Cor}
Let $\alpha\ge-1/2$, $\beta\in(-\frac12,\frac12)$, $4/3 < p < 4$ with
$\alpha+\beta\ge-1/2$,
\begin{gather*}
   \max\left\{ -\frac14, -\frac14 - \frac\beta2, - \frac{\alpha+1}2 - \beta
       \right\}
   < (\alpha+1) \left(\frac12 - \frac1p \right) <
   \min\left\{ \frac14, \frac14 - \frac\beta2 \right\},
   \\
   \max\left\{
   -\frac14, -\frac14 + \frac\beta2, - \frac{\alpha+1}2 + \frac\beta2
       \right\}
   < (\alpha+\beta+1) \left(\frac12 - \frac1p \right) <
   \min\left\{ \frac14, \frac14 + \frac\beta2 \right\}.
\end{gather*}
If $p<2$, assume further $-\frac14 + \frac\beta2 <
(\alpha+1) \left(\frac12 - \frac1p \right)$ and $-\frac14 - \frac\beta2 <
(\alpha+\beta+1) \left(\frac12 - \frac1p \right)$.

Then, $E_{p,\alpha} \cap L^p(x^{\alpha+\beta}) = E_{p,\alpha+\beta} \cap
L^p(x^\alpha)$.
\end{Cor}

\begin{proof}
The inclusion ``$\subseteq$'' is clear by Corollary~\ref{cor:EpaEpa+b}.
The inclusion ``$\supseteq$'' follows also by Corollary~\ref{cor:EpaEpa+b} with
$\alpha+\beta$ instead of $\alpha$, and $\beta$ instead of $-\beta$.
\end{proof}

Theorem~\ref{th:Bp=Ep} gives different basis for $E_{p,\alpha}$ for different
values of $\beta$. It seems interesting to obtain the expressions for the
change of basis between $\{j_n^\alpha\}_{n=0}^\infty$ and
$\{j_n^{\alpha+\beta}\}_{n=0}^\infty$ in $E_{p,\alpha}$.

\begin{Cor}
Let $\alpha\ge-1/2$, $-1/2 < \beta < 1$, $4/3 < p < 4$ with
\begin{equation*}
   \max\left\{ -\frac14, -\frac14 - \frac\beta2, - \frac{\alpha+1}2 - \beta
       \right\}
   < (\alpha+1) \left(\frac12 - \frac1p \right) <
   \min\left\{ \frac14, \frac14 - \frac\beta2 \right\}.
\end{equation*}
If $p<2$, assume further $-\frac14 + \frac\beta2 < (\alpha+1) \left(\frac12 -
\frac1p \right)$.
Then, both $\{j_n^\alpha\}_{n=0}^\infty$ and
$\{j_n^{\alpha+\beta}\}_{n=0}^\infty$ are bases of the space $E_{p,\alpha}$
and the change of basis is given by
\begin{equation*}
    j_n^\alpha = \sum_{k=n}^\infty a_{n,k} j_k^{\alpha+\beta},
    \qquad
    j_n^{\alpha+\beta} = \sum_{k=n}^\infty b_{n,k} j_k^\alpha,
\end{equation*}
where
\begin{equation}
\label{ec:ankbnk}
\begin{split}
    a_{n,k} &= \frac{2^\beta \, \sqrt{\alpha + 2n + 1} \,
        \sqrt{\alpha+\beta+2k+1} \,
        \Gamma(1-\beta) \, \Gamma(\alpha+\beta+k+n+1)}
       {\Gamma(1+k-n) \, \Gamma(1-\beta-k+n)\,
        \Gamma(\alpha+k+n+2)},
    \\
    b_{n,k} &= \frac{ 2^{-\beta} \, \sqrt{\alpha+\beta+2n+1} \,
        \sqrt{\alpha+2k+1} \,
        \Gamma(1+\beta) \, \Gamma(\alpha+k+n+1) }
      { \Gamma(1+k-n) \, \Gamma(1+\beta-k+n)\,
        \Gamma(\alpha+\beta+k+n+2) }.
\end{split}
\end{equation}
\end{Cor}

\begin{proof}
By Theorem~\ref{th:Bp=Ep}, $B_{p,\alpha,0} = E_{p,\alpha}$, so that
$j_n^\alpha \in E_{p,\alpha}$. Also, by Theorem~\ref{th:Bp=Ep}, $E_{p,\alpha}
= B_{p,\alpha,\beta}$. Then, Corollary~\ref{cor:conv} gives $S_N j_n^\alpha
\to j_n^\alpha$ in $L^p(x^\alpha)$ as $N\to\infty$, that is,
\begin{equation*}
   j_n^\alpha = \sum_{k=0}^\infty a_{n,k} j_k^{\alpha+\beta}
\end{equation*}
in $L^p(x^\alpha)$, where
\begin{equation*}
   a_{n,k} =
   \int_0^\infty j_n^\alpha j_k^{\alpha+\beta} x^{\alpha+\beta} \, dx.
\end{equation*}
In a similar way, $j_n^{\alpha+\beta} \in B_{p,\alpha,\beta} = B_{p,\alpha,0}$
and $j_n^{\alpha+\beta} = \sum_{k=0}^\infty b_{n,k} j_k^\alpha$ in
$L^p(x^\alpha)$, where
\begin{equation*}
   b_{n,k} =
   \int_0^\infty j_n^{\alpha+\beta} j_k^\alpha x^\alpha \, dx.
\end{equation*}
Finally, \cite[Ch.~XIII, 13.41 (2), p.~403]{Wat} gives $a_{n,k} = b_{n,k} = 0$
for $k<n$ and~(\ref{ec:ankbnk}) for $k \geq n$.
\end{proof}

Similar expressions for the change of basis between different bases
$\{j_n^{\alpha+\beta}\}_{n=0}^\infty$ and
$\{j_n^{\alpha+\beta'}\}_{n=0}^\infty$ in $E_{p,\alpha}$ can be obtained.
Details are left to the reader.

%%%%%%%%%%%%%%%%%%%%%%%%%%%%%%%%%%%%%%%%%%%%%%%%%%%%%%%%%%%%%%%%%%%%
\section{Proof of Theorem~\ref{th:boundSn}}
\label{sec:proofboundSn}
%%%%%%%%%%%%%%%%%%%%%%%%%%%%%%%%%%%%%%%%%%%%%%%%%%%%%%%%%%%%%%%%%%%%

\subsection{Necessary conditions}

Let us begin by showing that conditions~(\ref{ec:range})
are necessary for the uniform boundedness in Theorem~\ref{th:boundSn}.
Assume $S_n$ is uniformly bounded. Then the operators given by
\begin{equation*}
   T_n(f,x) = S_n(f,x) - S_{n-1}(f,x) =
   j_n^{\alpha+\beta}(x) \int_0^\infty f(t) j^{\alpha+\beta}_n(t)
   t^{\alpha+\beta} \, dt
\end{equation*}
are uniformly bounded as well, i.e.
\begin{equation*}
    \lpnorm{T_nf}p{x^\alpha}
    = \left| \int_0^\infty f(t) j_n^{\alpha+\beta}(t) t^{\alpha+\beta} \, dt
      \right| \lpnorm{j_n^{\alpha+\beta}}p{x^\alpha}
    \le C \lpnorm{f}p{x^\alpha}
\end{equation*}
with a constant $C$ independent of $n$ and $f$.
By duality, this means
\begin{equation*}
    \lpnorm{t^\beta j_n^{\alpha+\beta}}q{t^\alpha}
    \lpnorm{j_n^{\alpha+\beta}}p{x^\alpha}
    \le C,
\end{equation*}
where $1/p+1/q=1$. Taking $n$ fixed (it suffices $n=0$) and applying the first
part of Lemma~\ref{lem:normajn} gives~(\ref{ec:range}). Now,
provided~(\ref{ec:range}) holds, the norm estimates of Lemma~\ref{lem:normajn}
give
\begin{equation*}
    \lpnorm{j_n^{\alpha+\beta}}{p}{x^\alpha} \sim
    \begin{cases}
    n^{-(\alpha+\beta+1)+2(\alpha+1)/p}, & \text{if $p<4$,} \\
    n^{-(\alpha+2\beta+1)/2}(\log n)^{1/4}, & \text{if $p=4$,} \\
    n^{-(5/6+\alpha+\beta)+(6\alpha+4)/(3p)}, & \text{if $p>4$,}
    \end{cases}
\end{equation*}
and
\begin{equation*}
    \lpnorm{x^\beta j_n^{\alpha+\beta}}{q}{x^\alpha} \sim
    \begin{cases}
    n^{-(\alpha-\beta+1)+2(\alpha+1)/q}, & \text{if $q<4$ (i.e., $p>4/3$)}, \\
    n^{-(\alpha-2\beta+1)/2}(\log n)^{1/4}, & \text{if $q=4$ (i.e., $p=4/3$)},
      \\
    n^{-(5/6+\alpha-\beta)+(6\alpha+4)/(3q)}, & \text{if $q>4$ (i.e.,
      $p<4/3$)}.
    \end{cases}
\end{equation*}
This implies $4/3<p<4$.

%--------------------

\subsection{Sufficient conditions}

Let us now assume $\frac43<p<4$ and~(\ref{ec:range}), and prove the uniform
boundedness of the partial sum operators $S_n$. They can be written as
\begin{equation*}
   S_n(f,x) = \int_0^\infty {f(t) K_n(x,t) t^{\alpha+\beta}\,dt}, \quad
   \text{where \quad}
   K_n(x,t) = \sum_{k=0}^n j_k^{\alpha+\beta}(x) j_k^{\alpha+\beta}(t).
\end{equation*}
The next lemma gives a suitable decomposition of the kernel $K_n(x,t)$
associated to $S_n$. For a similar formula, with a different proof,
see~\cite{Var}.

\begin{Lem}
\label{lem:decomp}
Let $n \in \mathbb{N}$ and $\lambda > -1$. Then
\begin{multline*}
   \sum_{k=0}^n 2 (\lambda+2k+1) J_{\lambda+2k+1}(x) J_{\lambda+2k+1}(t)
   \\
   = \frac{xt}{x^2-t^2} \, [
   x J_{\lambda+1}(x) J_\lambda(t) - t J_\lambda(x) J_{\lambda+1}(t)
   \\
   + x J'_{\lambda+2n+2}(x) J_{\lambda+2n+2}(t)
   - t J_{\lambda+2n+2}(x) J'_{\lambda+2n+2}(t)
   ].
\end{multline*}
\end{Lem}

\begin{proof}
Using the equality $J_{\nu-1}(z) + J_{\nu+1}(z) = \frac{2 \nu}z J_\nu(z)$
(see~\cite[Ch.~III, 3.2, p.~45]{Wat}) to express $J_{\mu-1}$ and $J_{\mu+2}$ in
terms of $J_\mu$ and $J_{\mu+1}$ yields the formula
\begin{gather*}
  \frac{xt}{x^2-t^2} \left[
   x J_\mu(x) J_{\mu-1}(t) - t J_{\mu-1}(x) J_\mu(t)
   - x J_{\mu+2}(x) J_{\mu+1}(t)
   + t J_{\mu+1}(x) J_{\mu+2}(t)
   \right]
   \\
   = 2 \mu J_\mu(x) J_\mu(t).
\end{gather*}
This now gives
\begin{multline*}
   \sum_{k=0}^n
   2 (\lambda+2k+1) J_{\lambda+2k+1}(x) J_{\lambda+2k+1}(t)
   \\
   = \frac{xt}{x^2-t^2} \, [
   x J_{\lambda+1}(x) J_\lambda(t) - t J_\lambda(x) J_{\lambda+1}(t)
   \\
   - x J_{\lambda+2n+3}(x) J_{\lambda+2n+2}(t)
   + t J_{\lambda+2n+2}(x) J_{\lambda+2n+3}(t)
   ].
\end{multline*}
Finally, use the formula $z J_{\nu+1}(z) = \nu J_\nu(z) - z J'_\nu(z)$
(see~\cite[Ch.~III, 3.2, p.~45]{Wat}) to take out $J_{\lambda+2n+3}$.
\end{proof}

{}From the definition, we have
\begin{multline*}
   S_n(f,x)
   = x^{-\alpha/2-\beta/2-1/2} \int_0^\infty
   \bigg( \sum_{k=0}^n (\alpha+\beta+2k+1)
   \\
   \times J_{\alpha+\beta+2k+1}(\sqrt x)
   J_{\alpha+\beta+2k+1}(\sqrt t)
   \bigg) t^{\alpha/2+\beta/2-1/2} f(t) \, dt
\end{multline*}
so that Lemma~\ref{lem:decomp} with $\lambda=\alpha+\beta$ leads to
\begin{equation*}
   S_n f = W_1 f - W_2 f + W_{3,n} f - W_{4,n} f,
\end{equation*}
where
\begin{align*}
   W_1(f,x) &= \tfrac12 x^{-\alpha/2-\beta/2+1/2} J_{\alpha+\beta+1}(x^{1/2})
   \, H\big( t^{\alpha/2+\beta/2} J_{\alpha+\beta}(t^{1/2}) f(t), x \big),
   \\
   W_2(f,x) &= \tfrac12 x^{-\alpha/2-\beta/2} J_{\alpha+\beta}(x^{1/2})
   \, H\big(t^{\alpha/2+\beta/2+1/2} J_{\alpha+\beta+1}(t^{1/2}) f(t), x \big),
   \\
   W_{3,n}(f,x) &= \tfrac12 x^{-\alpha/2-\beta/2+1/2} J_\nu'(x^{1/2})
   \, H\big(t^{\alpha/2+\beta/2} J_\nu(t^{1/2}) f(t), x\big),
   \\
   W_{4,n}(f,x) &= \tfrac12 x^{-\alpha/2-\beta/2} J_\nu(x^{1/2})
   \, H\big(t^{\alpha/2+\beta/2+1/2} J_\nu'(t^{1/2}) f(t), x\big)
\end{align*}
and $\nu = \lambda+2n+2 = \alpha+\beta+2n+2$. Here, $H$ denotes the
Hilbert transform on $(0,\infty)$, which is defined by
\begin{equation*}
    H(g,x) = \int_0^\infty \frac{g(t)}{x-t}\,dt
\end{equation*}
(the integral must be considered as a principal value).

Thus, we can conclude that the partial sum operators $S_n$ are uniformly
bounded if we can prove that the operators $W_1$, $W_2$ are bounded and the
operators $W_{3,n}$, $W_{4,n}$ are uniformly bounded for $n \ge 0$. We will
use good estimates for the Bessel functions and the $A_p$ theory of weights to
prove the boundedness of the Hilbert transform.

Let us start with the bounds for the Bessel functions and their derivatives.

{}From the estimates~(\ref{ec:bessel0}) and~(\ref{ec:besselinfty}) it
follows that, for $\mu>-1$,
\begin{equation}
\label{ec:cota01}
    |J_\mu(x)| \le C_\mu x^\mu,
    \quad x \in (0,1],
\end{equation}
\begin{equation}
\label{ec:cota1infty}
    |J_\mu(x)| \le C_\mu x^{-1/2},
    \quad x \in [1,\infty),
\end{equation}
with a $C_\mu$ constant depending on $\mu$.

Moreover, we will need bounds for the Bessel functions
$J_{\alpha+\beta+2n+2}$ (and their derivatives) with constants independent of
$n$. So, we will make use of the bounds
\begin{equation}
\label{ec:jnu}
   |J_\nu(x)| \le C x^{-1/4} \left( |x - \nu| + \nu^{1/3} \right)^{-1/4},
   \quad x \in (0,\infty),
\end{equation}
\begin{equation}
\label{ec:jnu'}
   |J_\nu'(x)| \le C x^{-3/4} \left( |x - \nu| + \nu^{1/3} \right)^{1/4},
   \quad x \in (0,\infty),
\end{equation}
with some constant $C$ independent of~$\nu$. They follow from those
of~\cite{Bar-Cor,Cord}, for instance, and were already used in~\cite{Var}.

%--------------------

\subsection{Some results on Hilbert transforms and $A_p$ theory}

To analyze the boundedness of the Hilbert transform, some notation and
previous results will be necessary. As usual, for $1 < p < \infty$ we
write $q = p/(p-1)$, i.e., $1/p + 1/q = 1$. A weight is a nonnegative
Lebesgue-measurable function on $(0,\infty)$. The class $A_p(0,\infty)$
[$A_p$, for short] consists of those weights $w$ such that, for
every subinterval $I \subseteq (0, \infty)$,
\begin{equation*}
   \left(\frac1{|I|} \int_I w \right)
   \left( \frac1{|I|} \int_I w^{-q/p} \right)^{p/q} \le C,
\end{equation*}
where $C$ is a positive constant independent of $I$, and $|I|$ denotes the
length of $I$. The $A_p$ constant of $w$ is the least constant $C$
verifying this inequality and will be denoted by $A_p(w)$. We refer the reader
to~\cite{GC-RF} for further details on $A_p$ classes.

Fix $1 < p < \infty$; then the Hilbert transform $H$ is a bounded linear
operator on $L^p(w)$, for any weight $w \in A_p$. The norm of $H \colon L^p(w)
\longrightarrow L^p(w)$ and the $A_p$ constant of $w$ depend only on
each other, in the sense that given some constant $C$ which verifies the $A_p$
condition for $w$, another constant $C_1$ depending only on $C$ can be chosen
so that $\| H\| \le C_1$, and viceversa. Therefore, for a sequence $\{w_n\}_{n
\in \mathbb{N}}$ uniformly in $A_p$, i.e., with some constant $C$ verifying
the $A_p$ condition for every $w_n$, the Hilbert transform is uniformly
bounded on $L^p(w_n)$, $n \in \mathbb{N}$.

Let us see some auxiliary results related with $A_p$ weights:

\begin{Lem}
\label{lem:Ap1}
Let $u$, $v$, $w$ be weights on $(0, \infty)$ and $\gamma$ be a positive
constant. Then
\begin{itemize}
\item[(a)] $w(x) \in A_p$ if and only if $w(\gamma x) \in A_p$;
both weights have the same $A_p$ constant.
\item[(b)] $w \in A_p$ if and only if $\gamma w \in A_p$; both weights
have also the same $A_p$ constant.
\item[(c)] If $u, v \in A_p$, then $u + v \in A_p$ and
$A_p(u + v) \le A_p(u) + A_p(v)$.
\item[(d)] If $u, v \in A_p$ and $1/w = 1/u + 1/v$, then $w \in A_p$
and $A_p(w) \le C [A_p(u) + A_p(v)]$.
\end{itemize}
\end{Lem}

\begin{proof}
Parts (a) and (b) are trivial. Part (c) follows easily from
the inequality
\begin{equation*}
   \left(
   \frac1{|I|} \int_I (u + v)^{-{q/p}} \right)^{p/q}
   \le
   \min\left\{
   \left( \frac1{|I|} \int_I u^{-{q/p}} \right)^{p/q},
   \left( \frac1{|I|} \int_I v^{-{q/p}} \right)^{p/q}
   \right\}.
\end{equation*}
Part (d) is a consequence of (c) and the fact that $u \in A_p \iff
u^{-q/p} \in A_{q}$, with $A_{q}(u^{-q/p}) = [A_p(u)]^{q/p}$.
\end{proof}

The proof of the next lemma is not difficult, but cumbersome, so we omit it.
For the weight in~(c), observe that $x^r |x^{1/2} - 1|^s \sim x^r$ near
$0$, $x^r |x^{1/2} - 1|^s \sim |x-1|^s$ near $1$ and $x^r |x^{1/2} - 1|^s \sim
x^{r+s/2}$ near $\infty$, whence the three conditions follow.

\begin{Lem}
\label{lem:Ap2}
Let $r, s \in \mathbb{R}$. We have
\begin{itemize}
\item[(a)] $x^r \in A_p \iff -1 < r < p-1$.
\item[(b)] Set $\Phi(x) = x^r$ if $x \in (0,1)$ and $\Phi(x) = x^s$ if $x
\in (1,\infty)$. Then, $\Phi \in A_p$ if and only if $-1 < r < p-1$ and $-1 <
s < p-1$.
\item[(c)] $x^r |x^{1/2} - 1|^s \in A_p \iff -1 < r < p-1, \ -1 < s < p-1$ and
$-1 < r + s/2 < p-1$.
\end{itemize}
\end{Lem}

To simplify the notation, in the rest of this section we write
$\lambda=\alpha+\beta$ and $\nu=\alpha+\beta+2n+2$.

%--------------------

\subsection{Boundedness of the operators $W_1$ and $W_2$.}

{}From the definition, it follows that
\begin{equation*}
   \lpnorm{W_1f}p{x^\alpha} \le C \lpnorm{f}p{x^\alpha}
\end{equation*}
if and only if
\begin{equation*}
   \lpnorm{Hg}p{x^{\alpha - \lambda p/2 + p/2} |J_{\lambda+1}(x^{1/2})|^p}
   \le C \lpnorm{g}p{x^{\alpha - \lambda p/2} |J_\lambda(x^{1/2})|^{-p}}.
\end{equation*}
Proving that there is a weight $\Phi \in A_p$ with
\begin{equation}
\label{ec:cotaPhi}
   C x^{\alpha - \lambda p/2 + p/2} |J_{\lambda+1}(x^{1/2})|^p
   \le \Phi(x) \le
   C_1 x^{\alpha - \lambda p/2} |J_\lambda(x^{1/2})|^{-p}
\end{equation}
will be enough. According to the bounds~(\ref{ec:cota01})
and~(\ref{ec:cota1infty}), we have
\begin{align*}
   x^{\alpha - \lambda p/2 + p/2} |J_{\lambda+1}(x^{1/2})|^p
   &\le
   \begin{cases}
   C x^{\alpha + p}, &{\text{ if }} x \in (0,1), \\
   C x^{\alpha - \lambda p/2 + p/4}, &{\text{ if }} x \in (1,\infty),
   \end{cases}   
   \\
   x^{\alpha - \lambda p/2} |J_\lambda(x^{1/2})|^{-p}
   &\ge
   \begin{cases}
   C x^{\alpha - \lambda p}, &{\text{ if }} x \in (0,1), \\
   C x^{\alpha - \lambda p/2 + p/4}, &{\text{ if }} x \in (1,\infty).
   \end{cases}
\end{align*}
Let us try
\begin{equation*}
   \Phi(x) = \begin{cases}
   x^r, &{\text{ if }} x \in (0,1), \\
   x^{\alpha - \lambda p/2 + p/4},
   &{\text{ if }} x \in (1,\infty).
   \end{cases}
\end{equation*}
By (b) in Lemma~\ref{lem:Ap2}, conditions~(\ref{ec:cotaPhi}) and $\Phi \in A_p$
will hold if
\begin{equation}
\label{ec:ineqW1}
   \begin{cases}
   \alpha - \lambda p \le r \le \alpha + p, \\
   -1 < r < p - 1, \\
   -1 < \alpha - \lambda p/2 + p/4 < p - 1.
   \end{cases}
\end{equation}
The third line is equivalent to
\begin{equation*}
    \frac{-1}4 < (\alpha+1) \left( \frac12 - \frac1p \right) + \frac\beta2
    < \frac34,
\end{equation*}
which follows from~(\ref{ec:range}).
For the inequalities in~(\ref{ec:ineqW1}) involving $r$, it suffices to show
that
\begin{equation*}
   \max\{-1, \alpha - \lambda p\} < \min\{p - 1, \alpha + p\}.
\end{equation*}
This follows from $\alpha > -1$, $\alpha+\beta > -1$, $1 < p < \infty$,
and~(\ref{ec:range}), as well. The case of $W_2$ is entirely similar.

%--------------------

\subsection{Uniform boundedness of the operators $W_{3,n}$.}

Here,
\begin{equation*}
   \lpnorm{W_{3,n}f}p{x^\alpha} \le C \lpnorm{f}p{x^\alpha}
\end{equation*}
if and only if
\begin{equation*}
   \lpnorm{Hg}p{x^{\alpha - \lambda p/2 + {p/2}} |J_\nu'(x^{1/2})|^p}
   \le C \lpnorm{g}p{x^{\alpha - \lambda p/2} |J_\nu(x^{1/2})|^{-p}}.
\end{equation*}
{}From the bounds~(\ref{ec:jnu'}) and~(\ref{ec:jnu}),
\begin{align*}
   x^{\alpha - \lambda p/2 + {p/2}} |J_\nu'(x^{1/2})|^p
   &\le C x^{\alpha - \lambda p/2 + p/8}
   \left[|x^{1/2} - \nu| + \nu^{1/3} \right]^{p/4},
   \\
   x^{\alpha - \lambda p/2} |J_\nu(x^{1/2})|^{-p}
   &\ge C x^{\alpha - \lambda p/2 + p/8}
   \left[|x^{1/2} - \nu| + \nu^{1/3} \right]^{p/4}.
\end{align*}
It will suffice to prove that $\varphi_\nu \in A_p$ uniformly in $n$ (recall
that $\nu = \alpha + \beta + 2n + 2$), with
\begin{equation*}
   \varphi_\nu(x) = x^{\alpha - \lambda p/2 + p/8}
   \left[|x^{1/2} - \nu| + \nu^{1/3} \right]^{p/4}.
\end{equation*}
{}From Lemma~\ref{lem:Ap1}, we have
\begin{align*}
   \varphi_\nu(x) \in A_p {\text{ unif.}}
   &\iff \varphi_\nu(\nu^2 x) \in A_p {\text{ unif.}}
   \\
   &\iff x^{\alpha - \lambda p/2 + p/8}
   \left[ |x^{1/2} - 1| + \nu^{-2/3} \right]^{p/4} \in A_p {\text{ unif.}}
   \\
   &\iff x^{\alpha - \lambda p/2 + p/8} |x^{1/2} - 1|^{p/4}
   + \nu^{-p/6} x^{\alpha - \lambda p/2 + p/8}
   \in A_p {\text{ unif.}},
\end{align*}
where the last equivalence follows from
\begin{equation*}
   \left[ |x^{1/2} - 1| + \nu^{-2/3} \right]^{p/4}
   \sim |x^{1/2} - 1|^{p/4} + \nu^{-p/6}.
\end{equation*}
Now, Lemmas~\ref{lem:Ap1} and~\ref{lem:Ap2} prove that those weights belong to
$A_p$ uniformly.

%--------------------

\subsection{Uniform boundedness of the operators $W_{4,n}$.}

Finally,
\begin{equation*}
   \lpnorm{W_{4,n}f}p{x^\alpha} \le C \lpnorm{f}p{x^\alpha}
\end{equation*}
if and only if
\begin{equation*}
   \lpnorm{Hg}p{x^{\alpha - \lambda p/2} |J_\nu(x^{1/2})|^p}
    \le C \lpnorm{g}p{x^{\alpha - \lambda p/2 - p/2} |J_\nu'(x^{1/2})|^{-p}}.
\end{equation*}
Also,
\begin{align*}
   x^{\alpha - \lambda p/2} |J_\nu(x^{1/2})|^p
   &\le C x^{\alpha - \lambda p/2 - p/8}
   \left[|x^{1/2} - \nu| + \nu^{1/3} \right]^{-p/4},
   \\
   x^{\alpha - \lambda p/2 - {p/2}} |J_\nu'(x^{1/2})|^{-p}
   &\ge C x^{\alpha - \lambda p/2 - p/8}
   \left[|x^{1/2} - \nu| + \nu^{1/3} \right]^{-p/4},
\end{align*}
so let us put
\begin{equation*}
   \psi_\nu(x) = x^{\alpha - \lambda p/2 - p/8}
   \left[ |x^{1/2} - \nu| + \nu^{1/3} \right]^{-p/4}
\end{equation*}
and show that $\psi_\nu \in A_p$ uniformly in $n$. Indeed,
\begin{align*}
   \psi_\nu(x) \in A_p {\text{ unif.}}
   &\iff \psi_\nu(\nu^2 x) \in A_p {\text{ unif.}}
   \\
   &\iff x^{\alpha - \lambda p/2 - p/8}
   \left[ |x^{1/2} - 1| + \nu^{-2/3} \right]^{-p/4} \in A_p {\text{ unif.}}
\end{align*}
and
\begin{multline*}
   \Big(x^{\alpha - \lambda p/2 - p/8}
   \left[ |x^{1/2} - 1| + \nu^{-2/3} \right]^{-p/4}
   \Big)^{-1}
   \\
   \sim x^{-\alpha + \lambda p/2 + p/8}
   \left[|x^{1/2} - 1|^{p/4} + \nu^{-p/6} \right]
   \\
   = \left[x^{\alpha - \lambda p/2 - p/8}
   |x^{1/2} - 1|^{-p/4} \right]^{-1}
   + \left[ \nu^{p/6} x^{\alpha - \lambda p/2 - p/8}
   \right]^{-1},
\end{multline*}
so that Lemmas~\ref{lem:Ap1} and~\ref{lem:Ap2} lead to the desired conclusion.
The proof of Theorem~\ref{th:boundSn} is now complete.

%%%%%%%%%%%%%%%%%%%%%%%%%%%%%%%%%%%%%%%%%%%%%%%%%%%%%%%%%%%%%%%%%%%%
\section{Proof of Theorem~\ref{th:convae}}
\label{sec:proofconvae}
%%%%%%%%%%%%%%%%%%%%%%%%%%%%%%%%%%%%%%%%%%%%%%%%%%%%%%%%%%%%%%%%%%%%

We only need to prove that $S_n(f,x)$ converges to some $g(x)$ almost
everywhere. This, together with Corollary~\ref{cor:conv}, gives $g = f$ almost
everywhere. Now, recall that $S_n f = \sum_{k=0}^n c_k(f) j_k^{\alpha+\beta}$,
where
\begin{equation*}
   c_k(f) = \int_0^\infty f(t) j_k^{\alpha+\beta}(t) t^{\alpha+\beta} \, dt.
\end{equation*}
It follows from Lemma~\ref{lem:normajn} that
$x^\beta j_n^{\alpha+\beta} \in L^q(x^\alpha)$; moreover,
$\lpnorm{x^\beta j_n^{\alpha+\beta}}{q}{x^\alpha} \le C n^\delta$
for some constant $\delta = \delta(p,\alpha,\beta)$. Thus,
\begin{equation*}
   |c_n(f)|
   \le \lpnorm{f}p{x^\alpha}
       \lpnorm{x^\beta j_n^{\alpha+\beta}}q{x^\alpha}
   \le C \lpnorm{f}p{x^\alpha}\ n^\delta.
\end{equation*}
Now, according to~\cite[Ch.~III, 3.31 (1), p.~49]{Wat} we have
\begin{equation*}
   |J_\nu(x)| \leq \frac{2^{-\nu} x^\nu}{\Gamma(\nu+1)},
   \qquad \nu > -1/2.
\end{equation*}
Therefore,
\begin{align*}
   |j_n^{\alpha+\beta}(x)|
   &= \sqrt{\alpha+\beta+2n+1}\ |J_{\alpha+\beta+2n+1}(\sqrt{x})|
   \ x^{-\frac{\alpha+\beta+1}{2}} \\
   &\leq \frac{\sqrt{\alpha+\beta+2n+1}\ 2^{-(\alpha+\beta+2n+1)} x^n}
             {\Gamma(\alpha+\beta+2n+2)},
\end{align*}
so that
\begin{equation*}
   |c_n(f) j_n^{\alpha+\beta}(x)| \leq C \lpnorm{f}p{x^\alpha}
   \ \frac{n^{\delta+1/2} (x/4)^n}{\Gamma(\alpha+\beta+2n+2)}
\end{equation*}
and the series $\sum_{n=0}^\infty c_n(f) j_n^{\alpha+\beta}(x)$ converges.

%%%%%%%%%%%%%%%%%%%%%%%%%%%%%%%%%%%%%%%%%%%%%%%%%%%%%%%%%%%%%

\begin{Ackn} The authors thank the referees for their valuable
comments, which helped us to improve this paper in its final version.
\end{Ackn}

%%%%%%%%%%%%%%%%%%%%%%%%%%%%%%%%%%%%%%%%%%%%%%%%%%%%%%%%%%%%%

\bibliographystyle{amsplain}

\begin{thebibliography}{99}

\bibitem{Bar-Cor} J. A. Barcel\'o and A. C\'ordoba,
{Band-limited functions: $L^p$-convergence},
\textit{Trans. Amer. Math. Soc.}
\textbf{313} (1989), 655--669.

\bibitem{Cord} A. C\'ordoba,
{The disc multiplier},
\textit{Duke Math. J.}
\textbf{58} (1989), 21--27.

\bibitem{Dur} A. J. Dur\'an,
{On Hankel transform},
\textit{Proc. Amer. Math. Soc.}
\textbf{110} (1990), 417--424.

\bibitem{E-HTF} A. Erd\'elyi, W. Magnus, F. Oberhettinger and F. Tricomi,
{``Higher Transcendental Functions,''} Vol.~II,
McGraw-Hill, New York, 1953.

\bibitem{E-TIT} A. Erd\'elyi, W. Magnus, F. Oberhettinger and F. Tricomi,
{``Tables of Integral Transforms,''} Vol.~II,
McGraw-Hill, New York, 1954.

\bibitem{GC-RF} J. Garc\'{\i}a-Cuerva and J. L. Rubio de Francia,
{``Weighted Norm Inequalities and Related Topics,''}
Math. Studies \textbf{116},
North-Holland, Amsterdam, 1985.

\bibitem{Herz} C. S. Herz,
{On the mean inversion of Fourier and Hankel transforms},
\textit{Proc. Nat. Acad. Sci. U.S.A.}
\textbf{40} (1954), 996--999.

\bibitem{Me-So} J. M. M\'endez and M. M. Socas,
{A pair of generalized Hankel-Clifford transformations and their applications},
\textit{J. Math. Anal. Appl.}
\textbf{154} (1991), 543--557.

\bibitem{Muck} B. Muckenhoupt,
{Mean convergence of Jacobi series},
\textit{Proc. Amer. Math. Soc.}
\textbf{23} (1969), 306--310.

\bibitem{Pol-III} H. Pollard,
{The mean convergence of orthogonal series. III},
\textit{Duke Math. J.}
\textbf{16} (1949), 189--191.

\bibitem{Stem} K. Stempak,
{A uniform $L^p$ estimate of Bessel functions: a note on Guo's paper},
\textit{Proc. Amer. Math. Soc.} (to appear).

\bibitem{St-Tr} K. Stempak and W. Trebels,
{Hankel multipliers and transplanting operators},
\textit{Studia Math.}
\textbf{126} (1997), 51--66.

\bibitem{Sze} G. Szeg\H{o},
{``Orthogonal Polynomials''} (4th edition),
Amer. Math. Soc. Colloq. Publ. \textbf{23},
Amer. Math. Soc.,
Providence, RI, 1975.

\bibitem{Var} J. L. Varona,
{Fourier series of functions whose Hankel transform is supported on $[0,1]$},
\textit{Constr. Approx.}
\textbf{10} (1994), 65--75.

\bibitem{Wat} G. N. Watson,
{``A Treatise on the Theory of Bessel Functions''} (2nd edition),
Cambridge Univ. Press,
Cambridge, 1944.

\end{thebibliography}

\end{document}

%%%%%%%%%%%%%%%%%%%%%%%%%%%%%%%%%%%%%%%%%%%%%%%%%%%%%%%%%%%%%
% End of the paper
%%%%%%%%%%%%%%%%%%%%%%%%%%%%%%%%%%%%%%%%%%%%%%%%%%%%%%%%%%%%%